\documentclass[conference]{IEEEtran}
\IEEEoverridecommandlockouts
\usepackage{cite}
\usepackage{amsmath,amssymb,amsfonts}
\usepackage{graphicx}
\usepackage{booktabs}
\usepackage{textcomp}
\usepackage{xcolor}

\begin{document}

\title{Two-Sided Shrinkage Weights for Hybrid Recommendation\\
under Continuous Cold Start}

\author{\IEEEauthorblockN{Hoang Phuc}
\IEEEauthorblockA{\textit{Department of Software Engineering} \\
\textit{FPT University}\\
Ho Chi Minh City, Vietnam \\
dhphuc1406@gmail.com}
}

\maketitle

\begin{abstract}
A new hotel booking platform faces cold start on both sides at once:
travelers register without any history, and partner hotels join
without any booking. Collaborative filtering cannot score either, so
hybrid recommenders fall back on content and demographic signals.
Existing weighted hybrids, however, set the mixing weights either
once for the whole system or as a function of the user alone, which
still lets collaborative filtering dominate the ranking of an item
nobody has rated. We propose a fusion rule in which the collaborative
weight is the \emph{product} of two shrinkage factors, one over the
user's interaction count and one over the item's rating count, so
that collaborative filtering is trusted only when both sides carry
evidence. The remaining mass is split between content and
demographic components by a second shrinkage factor, giving a
three-weight rule with three interpretable constants. On MovieLens
100K with 20\% of items fully held out to simulate new items,
averaged over ten seeds, the proposed rule attains the lowest RMSE
(1.0175) among twelve methods, and places previously unrated items
in the top-10 at a rate of 0.0124, more than four times that of the
strongest published baseline and against exactly zero for every
single-paradigm and fixed-weight method. Gains on both measures are
significant under a Wilcoxon signed-rank test over seeds.
\end{abstract}

\begin{IEEEkeywords}
recommender systems, hybrid recommendation, cold start, shrinkage,
matrix factorization, hotel booking
\end{IEEEkeywords}

%=====================================================================
\section{Introduction}

% TODO: doan nay can so lieu that cua nen tang StayReco —
% bao nhieu partner, bao nhieu tinh, quy mo du kien.
% Hien dang de chung chung, sua lai khi co so.

Online booking platforms that launch with a new catalogue face a
structural difficulty: almost every cell of the user--item matrix is
empty, and it stays that way. New travelers register daily and new
partner hotels are verified continuously, so a substantial share of
every request involves at least one side with no interaction
history. Bernardi \textit{et al.} call this the \emph{continuous}
cold-start problem, in contrast to the one-off cold start assumed by
much of the literature \cite{bernardi2015}.

Collaborative filtering (CF) is the strongest signal once both the
user and the item are warm, but it cannot score an item that nobody
has interacted with \cite{schein2002}. Content-based filtering
handles new items immediately, since a hotel acquires a feature
vector as soon as its description is entered, but it ignores the
collective taste signal and tends to over-specialize
\cite{lops2011}. Demographic filtering requires nothing from the
user beyond a profile and therefore always applies, at the cost of
being the least accurate of the three \cite{pazzani1999}. A hybrid
recommender combines them \cite{burke2002}, but the combination rule
determines whether the hybrid inherits the strengths or the
weaknesses of its parts.

Fixed mixing weights are the common choice and the obvious weakness.
A single weight vector applied to every request necessarily
over-trusts CF for cold pairs and under-uses it for warm ones.
Melville \textit{et al.} report that a naive equal-weight blend of
collaborative and content predictions is in fact \emph{worse} than
pure CF on their data \cite{melville2002}, which is a direct argument
for making the weights depend on the evidence available.

Weights that adapt to the user's interaction count address half of
the problem. They do not address the other half: a long-standing
user browsing a newly listed hotel still receives a high
collaborative weight, applied to a latent factor that was never
trained. This paper argues that under two-sided cold start the
mixing weight must be a function of \emph{both} counts, and that the
natural form for each factor is the shrinkage ratio $n/(n+k)$ that
Koren uses to discount similarities estimated from few observations
\cite{koren2008}.

The contributions are:

\begin{enumerate}
\item A fusion rule whose collaborative weight is the product of a
user-side and an item-side shrinkage factor, so that a hotel with no
ratings receives zero collaborative weight regardless of how active
the user is.
\item A two-stage weighting scheme that applies the same shrinkage
form twice --- once to separate CF from the rest, once to divide the
rest between content and demographic signals --- yielding three
weights controlled by three interpretable constants.
\item An evaluation over ten seeds against eleven baselines,
including four recent hybrid methods reimplemented inside a single
harness, reported separately by cold-start group and tested with a
Wilcoxon signed-rank test.
\end{enumerate}

%=====================================================================
\section{Related Work}

\subsection{Hybrid Combination Strategies}

Burke surveys hybrid recommenders and classifies them into seven
families \cite{burke2002}. The \emph{weighted} hybrid, in which
``the scores (or votes) of several recommendation techniques are
combined together to produce a single recommendation,'' is the family
this work belongs to. Burke also states the weakness that motivates
adaptive weights: a fixed weighting implicitly assumes that the
relative value of each technique is uniform across the item space,
which is false precisely in the cold-start region.

Pazzani combines collaborative, content-based and demographic
components \cite{pazzani1999}, the same three paradigms used here,
although the combination is by consensus voting over each
recommender's output rather than by a weighted sum of scores.

\subsection{Weights That Depend on Evidence}

The P-Tango system of Claypool \textit{et al.} starts from equal
weights for the collaborative and content recommenders and adjusts
them per user as predictions are confirmed or disconfirmed
\cite{claypool1999}. Melville \textit{et al.} scale the confidence
placed in a content-based predictor by the number of items the
active user has rated, capped at a threshold of 50
\cite{melville2002}. Both adapt on the user axis only.

Within collaborative filtering itself, discounting estimates computed
from little data is standard. Koren shrinks an item--item
correlation by the factor $n_{ij}/(n_{ij}+\lambda)$, where $n_{ij}$
is the number of users who rated both items, on the stated ground
that ``similarities based on a greater user support are more
reliable'' \cite{koren2008}. The present work applies the same form
at the level of hybrid fusion weights rather than similarities, and
applies it to both axes.

\subsection{Recent Hybrid Recommenders}

Four recent methods are reimplemented as baselines. Tarek
\textit{et al.} combine item--item CF, content and a knowledge-based
fallback with fixed weights for hotel recommendation
\cite{tarek2022}. Do \textit{et al.} adjust the weights of an
SVD-plus-content hybrid according to validation performance
\cite{do2020}. Muhammad \textit{et al.} blend SVD and SVD++ with a
grid-searched mixing parameter, entirely within the collaborative
family \cite{ohwsf2025}. Çakır \textit{et al.} combine SVD++ and
LightGCN with a learned weight and switch on a demographic and
TF-IDF safety layer in cold-start situations; the switch is binary
rather than continuous \cite{slarec2026}.

\subsection{Evaluation of Cold Start}

Schein \textit{et al.} define cold-start evaluation protocols
\cite{schein2002}; Herlocker \textit{et al.} survey accuracy metrics
for collaborative filtering \cite{herlocker2004}. Krichene and Rendle
show that ranking against a sampled candidate set distorts
comparisons and recommend ranking over the full catalogue
\cite{krichene2020}, which this work follows. Ferrari Dacrema
\textit{et al.} document how often reported gains fail to reproduce
when baselines are tuned with the same effort as the proposed method
\cite{dacrema2019}; the protocol in Section~\ref{sec:setup} applies
one identical tuning procedure to all twelve methods.

%=====================================================================
\section{Problem Formulation}

Let $U$ be the set of users and $I$ the set of items. For a user
$u \in U$ the system returns an ordered list of $K$ items drawn from
the candidate set $C_u \subseteq I$ of items $u$ has not yet
consumed. Two evidence counts drive the method:

\begin{itemize}
\item $n_u$ --- the number of distinct items user $u$ has interacted
with in the training data;
\item $m_i$ --- the number of distinct users who have interacted with
item $i$ in the training data.
\end{itemize}

An item with $m_i = 0$ is a \emph{new item}: no collaborative model
has ever observed it. A user with small $n_u$ is a \emph{cold user}.
Both conditions occur simultaneously and continuously on a growing
platform, which is the setting this paper addresses.

In the target deployment a knowledge-based pre-filter is applied
before scoring, removing hotels that are unavailable for the
requested dates, outside the stated budget, or not yet approved.
These are hard business constraints rather than predictions, and they
are not part of the ranking model; they are not evaluated here
because the public dataset used in Section~\ref{sec:setup} carries no
availability or price-constraint information.

%=====================================================================
\section{Proposed Method}

\subsection{Component Models}

Three component models produce a score on the rating scale for every
candidate pair $(u,i)$.

\textbf{Collaborative (CF).} SVD++ \cite{koren2008} is trained on the
explicit rating matrix with $50$ latent factors, $30$ epochs,
learning rate $0.005$ and regularization $0.02$. It is the most
accurate component when both sides are warm and is undefined when
$m_i = 0$.

\textbf{Content-based (CB).} Items are represented by TF-IDF vectors
over their textual and categorical attributes. A user profile is the
interaction-weighted mean of the vectors of the items that user has
consumed; the predicted score is a baseline estimate plus a term
proportional to the cosine similarity between profile and item. The
component requires no other user's data, so it scores a new item
immediately.

\textbf{Demographic (DM).} The predicted score is the mean rating of
the user's demographic group for that item, shrunk toward the global
mean with a pseudo-count of $5$. It requires nothing from the user's
history and is therefore always available, at the lowest accuracy of
the three.

All three components are expressed on the same rating scale before
fusion, so that the weights carry their intended meaning. Their
empirical means over the test split are $3.49$, $3.54$ and $3.55$
respectively, against a data mean of $3.55$.

\subsection{Fusion Rule}

The predicted score is a convex combination of the three components:

\begin{equation}
\hat{s}_{u,i} \;=\; \alpha\, s^{\mathrm{CF}}_{u,i}
                 \;+\; \beta \, s^{\mathrm{CB}}_{u,i}
                 \;+\; \gamma\, s^{\mathrm{DM}}_{u,i},
\qquad
\alpha + \beta + \gamma = 1 .
\label{eq:fusion}
\end{equation}

\subsection{Two-Sided Shrinkage}

Define two trust factors, one per axis:

\begin{equation}
t^{\mathrm{user}}_u = \frac{n_u}{n_u + k}, \qquad
t^{\mathrm{item}}_i = \frac{m_i}{m_i + k'} .
\label{eq:trust}
\end{equation}

Each rises from $0$ toward $1$ as evidence accumulates and equals
$\tfrac{1}{2}$ when the count reaches the corresponding constant,
which gives $k$ and $k'$ a direct reading: the amount of evidence at
which that side is trusted half as much as a fully observed one. The
form follows Koren's shrinkage factor \cite{koren2008}.

The collaborative weight is their \emph{product}:

\begin{equation}
\alpha_{u,i} \;=\; t^{\mathrm{user}}_u \cdot t^{\mathrm{item}}_i .
\label{eq:alpha}
\end{equation}

The product is the operative choice. Because either factor vanishing
drives $\alpha$ to zero, an item with $m_i = 0$ receives no
collaborative weight however large $n_u$ is, and a user with
$n_u = 0$ receives none however well observed the item is. A sum or
a mean of the two factors would not have this property: a highly
active user would still carry collaborative weight onto an item the
model has never seen.

\subsection{Dividing the Remainder}

The mass not assigned to CF is divided between the content and
demographic components by a second application of the same form:

\begin{equation}
\pi_u = \frac{n_u}{n_u + k_c}, \qquad
\beta = (1-\alpha)\,\pi_u, \qquad
\gamma = (1-\alpha)\,(1 - \pi_u) .
\label{eq:beta}
\end{equation}

The content profile is itself built from the user's history, so it
improves with $n_u$; the demographic component does not depend on
the user's history at all. Within the remainder, therefore, the same
logic applies one level down: evidence shifts weight from the
unconditional fallback to the data-driven component. Three
constants, $k$, $k'$ and $k_c$, control the whole rule.

One exception is imposed. When $m_i = 0$ we set $\gamma = 0$ and
$\beta = 1 - \alpha$. For an item nobody has rated, the demographic
component has no item-specific information and returns the same
global-mean-dominated value for every such item; including it adds a
constant offset that cannot separate new items from one another and
dilutes the only signal that can, namely content similarity.
Section~\ref{sec:results} quantifies the effect: removing this
exception reduces the rate at which new items reach the top-10 by an
order of magnitude.

\subsection{Resulting Behaviour}

Table~\ref{tab:weights} shows the weights produced by
\eqref{eq:trust}--\eqref{eq:beta} under the constants selected in
Section~\ref{sec:setup}.

\begin{table}[t]
\caption{Fusion weights under the selected constants
($k=5$, $k'=100$, $k_c=10$). The item is well observed
($m_i = 200$) except in the last row.}
\label{tab:weights}
\centering
\begin{tabular}{lccc}
\toprule
Situation & $\alpha$ & $\beta$ & $\gamma$ \\
\midrule
New user ($n_u = 0$)        & 0.00 & 0.00 & 1.00 \\
$n_u = 10$                  & 0.44 & 0.28 & 0.28 \\
$n_u = 40$                  & 0.59 & 0.33 & 0.08 \\
Established user            & 0.65 & 0.33 & 0.02 \\
New item ($m_i = 0$)        & 0.00 & 1.00 & 0.00 \\
\bottomrule
\end{tabular}
\end{table}

The transition is continuous in both counts: an item with a single
rating receives a collaborative weight near zero rather than the full
weight a hard threshold would assign.

%=====================================================================
\section{Experimental Setup}
\label{sec:setup}

\subsection{Dataset}

% GHI CHU QUAN TRONG: day la diem yeu phai noi ro.
% MovieLens khong phai du lieu khach san.

The platform has no production interaction logs yet, so the method
is evaluated on MovieLens 100K \cite{harper2015}: $943$ users,
$1{,}682$ items and $100{,}000$ ratings. The dataset provides the
three inputs the method requires --- an interaction matrix, item
attributes usable as content features, and user demographic
attributes --- and is the standard benchmark for this class of
comparison. It is a proxy: hotel booking differs from film rating in
session structure, repeat-consumption behaviour and the role of
price and availability. Section~\ref{sec:threats} returns to this.

\subsection{Cold-Start Simulation}

Every MovieLens user has at least $20$ ratings, so cold start must be
induced. Two independent manipulations are applied per seed:

\begin{itemize}
\item \textbf{New items.} $20\%$ of items are selected at random and
\emph{all} of their ratings are moved to the test split, so these
items have $m_i = 0$ in training and no collaborative model has ever
observed them.
\item \textbf{Cold users.} $30\%$ of users are selected at random and
only their first $n$ ratings in chronological order are kept in
training, with $n$ drawn uniformly from $[0, 19]$.
\end{itemize}

The remaining data are split into training, validation and test
portions. Test rows are assigned to groups by the user's training
count: \emph{cold} ($n_u < 5$), \emph{cool}
($5 \le n_u < 20$), \emph{warm} ($n_u \ge 20$), plus a separate
\emph{new item} group.

\subsection{Protocol}

Every result is averaged over ten independent seeds, each
corresponding to a different random split.

All twelve methods are tuned by the identical procedure: a grid
search on the \emph{validation} split, selecting by RMSE. The test
split is not consulted during tuning. The grid for the proposed
method covers $k \in \{5,10,20,40,80,160\}$,
$k' \in \{5,20,50,100\}$ and $k_c \in \{1,3,10,30\}$, that is $96$
combinations; the selected values are $k=5$, $k'=100$, $k_c=10$.

\subsection{Metrics}

Rating accuracy is reported as RMSE and MAE. Ranking quality is
reported as Precision, Recall and NDCG at $K=10$, computed by ranking
over the \emph{full catalogue} of items the user has not rated,
following Krichene and Rendle \cite{krichene2020}. Ranking inside a
sampled candidate set of items the user happened to rate inflates all
methods and, critically, never contains an item with $m_i = 0$, so
it cannot measure item-side cold start at all.

Two beyond-accuracy measures are reported. Cov@10 is the fraction of
the catalogue that appears in any top-10 list. New@10 is the fraction
of top-10 slots occupied by items with $m_i = 0$; it is the direct
measure of whether a method can surface an item nobody has rated.

Differences between the proposed method and each baseline are tested
with a Wilcoxon signed-rank test over the ten per-seed values.

\subsection{Baselines}

Four single-paradigm references (Popularity, Demographic-only,
Content-only, CF-only), three fixed-weight hybrids over the same
three components ($\alpha \in \{0.3, 0.5, 0.7\}$), and four recent
hybrid methods \cite{tarek2022, do2020, ohwsf2025, slarec2026}
reimplemented inside this harness. The reimplementations are
adaptations, not reproductions: each original uses a different
dataset, split and metric, so their published numbers are not
comparable to ours and are not quoted. Every deviation required to
run them here is documented in the source.

%=====================================================================
\section{Results}
\label{sec:results}

\subsection{Overall}

Table~\ref{tab:overall} reports the mean over ten seeds.

\begin{table}[t]
\caption{Overall results, mean over ten seeds, $K=10$.
Best value per column in bold.}
\label{tab:overall}
\centering
\small
\begin{tabular}{lccccc}
\toprule
Method & RMSE & MAE & P@10 & NDCG@10 & New@10 \\
\midrule
Popularity          & 1.0823 & 0.8904 & 0.0320 & 0.0356 & 0.0000 \\
Demographic-only    & 1.0846 & 0.8920 & 0.0342 & 0.0406 & 0.0000 \\
Content-only        & 1.0241 & 0.8177 & 0.0342 & 0.0401 & 0.0018 \\
CF-only (SVD++)     & 1.0334 & 0.8312 & 0.0370 & 0.0450 & 0.0000 \\
\midrule
Fixed $\alpha = 0.3$ & 1.0187 & 0.8222 & 0.0372 & 0.0443 & 0.0000 \\
Fixed $\alpha = 0.5$ & 1.0196 & 0.8223 & 0.0382 & 0.0458 & 0.0000 \\
Fixed $\alpha = 0.7$ & 1.0232 & 0.8244 & 0.0383 & 0.0461 & 0.0000 \\
\midrule
Tarek+ \cite{tarek2022}    & 1.0231 & 0.8306 & 0.0363 & 0.0425 & 0.0000 \\
Do+ \cite{do2020}          & 1.0219 & 0.8197 & 0.0363 & 0.0432 & 0.0000 \\
OHWSF \cite{ohwsf2025}     & 1.0333 & 0.8312 & 0.0366 & 0.0446 & 0.0000 \\
SL-ARec \cite{slarec2026}  & 1.0181 & 0.8176 & \textbf{0.0412}
                           & \textbf{0.0505} & 0.0030 \\
\midrule
\textbf{Proposed} & \textbf{1.0175} & \textbf{0.8129} & 0.0403
                  & 0.0482 & \textbf{0.0124} \\
\bottomrule
\end{tabular}
\end{table}

The proposed rule attains the lowest RMSE and MAE of the twelve
methods. On New@10 it reaches $0.0124$ against $0.0030$ for the
strongest baseline, and against exactly $0.0000$ for every
single-paradigm method, every fixed-weight hybrid, and three of the
four published methods. Those nine methods never place an unrated
item in a top-10 list. On NDCG@10 the proposed rule is second,
behind SL-ARec \cite{slarec2026}.

\subsection{By Cold-Start Group}

Table~\ref{tab:groups} breaks NDCG@10 down by the user's training
count.

\begin{table}[t]
\caption{NDCG@10 by cold-start group, mean over ten seeds.}
\label{tab:groups}
\centering
\begin{tabular}{lccc}
\toprule
Method & Cold & Cool & Warm \\
       & ($n_u<5$) & ($5\!\le\! n_u\!<\!20$) & ($n_u \ge 20$) \\
\midrule
Popularity          & 0.0250 & 0.0209 & 0.0471 \\
Demographic-only    & 0.0320 & 0.0246 & 0.0525 \\
Content-only        & 0.0291 & 0.0217 & 0.0540 \\
CF-only (SVD++)     & 0.0228 & 0.0235 & 0.0625 \\
Fixed $\alpha=0.5$  & 0.0261 & 0.0238 & 0.0633 \\
SL-ARec \cite{slarec2026} & 0.0299 & 0.0233 & \textbf{0.0717} \\
\textbf{Proposed}   & \textbf{0.0324} & 0.0231 & 0.0673 \\
\bottomrule
\end{tabular}
\end{table}

The pattern matches the design. For cold users the proposed rule is
best of all twelve, because $\alpha$ is near zero there and the
ranking is carried by the content and demographic components; CF-only
is the weakest method in that column. For warm users CF-only improves
by a factor of nearly three over its cold-user figure, and the
proposed rule tracks it upward while remaining behind SL-ARec, whose
LightGCN component is stronger on dense data.

\subsection{Significance}

Table~\ref{tab:wilcoxon} reports the Wilcoxon signed-rank test over
the ten per-seed values.

\begin{table}[t]
\caption{Wilcoxon signed-rank test, proposed method against each
baseline, ten seeds. ``Wins'' counts seeds on which the proposed
method is better.}
\label{tab:wilcoxon}
\centering
\small
\begin{tabular}{lcccc}
\toprule
 & \multicolumn{2}{c}{RMSE} & \multicolumn{2}{c}{New@10} \\
\cmidrule(lr){2-3}\cmidrule(lr){4-5}
Baseline & Wins & $p$ & Wins & $p$ \\
\midrule
CF-only (SVD++)    & 10/10 & 0.002 & 10/10 & 0.002 \\
OHWSF \cite{ohwsf2025} & 10/10 & 0.002 & 10/10 & 0.002 \\
Tarek+ \cite{tarek2022} & 10/10 & 0.002 & 10/10 & 0.002 \\
Do+ \cite{do2020}  & 10/10 & 0.002 & 10/10 & 0.002 \\
Fixed $\alpha=0.3$ & 8/10  & 0.014 & 10/10 & 0.002 \\
SL-ARec \cite{slarec2026} & 6/10 & 0.322 & 10/10 & 0.002 \\
\bottomrule
\end{tabular}
\end{table}

On New@10 the advantage is significant against every baseline,
including SL-ARec. On RMSE it is significant against all baselines
except SL-ARec, where the difference over ten seeds is not
distinguishable from chance and no claim is made.

\subsection{Effect of the New-Item Exception}

Removing the $\gamma = 0$ rule of Section~IV-D, while leaving
everything else unchanged, lowers New@10 from $0.0124$ to $0.0011$,
a factor of eleven, and the difference is significant
($10/10$ seeds, $p = 0.002$). RMSE improves slightly in exchange,
from $1.0175$ to $1.0142$. The demographic component therefore does
carry a small amount of rating-scale accuracy on new items while
actively harming their ranking, which is consistent with it
returning a near-constant value for every such item.

%=====================================================================
\section{Limitations and Threats to Validity}
\label{sec:threats}

\textbf{Dataset.} The evaluation uses film ratings as a proxy for
hotel bookings. Explicit five-star ratings are not implicit booking
signals; films are consumed once and cheaply while hotel stays are
rare and expensive; and neither price constraints nor availability
exist in the data. The conclusions therefore concern the
\emph{weighting rule} and should not be read as measurements of the
deployed platform.

\textbf{Knowledge layer not evaluated.} The pre-filter described in
Section~III cannot be exercised on this dataset and is therefore
absent from all reported numbers.

\textbf{Grid boundary.} The selected $k' = 100$ is the largest value
in its grid, so the optimum may lie outside the searched range and
the reported configuration may be suboptimal.

\textbf{Single dataset.} All results come from one benchmark. The
magnitude of every difference reported here is small in absolute
terms, and generalization to other domains is untested.

\textbf{Reimplemented baselines.} The four published methods are
adaptations constrained to this harness. Their numbers here
characterize the \emph{idea} under common conditions, not the
original systems.

%=====================================================================
\section{Conclusion}

Under two-sided cold start the mixing weight of a hybrid recommender
should depend on the evidence available on both axes. Making the
collaborative weight the product of a user-side and an item-side
shrinkage factor enforces this directly: either side lacking data
removes the collaborative term entirely. Applying the same shrinkage
form a second time divides the remainder between a data-driven
content component and an unconditional demographic fallback, giving
a three-weight rule with three interpretable constants and no
thresholds.

On MovieLens 100K with new items simulated, the rule attains the
lowest rating error of twelve methods and raises the rate at which
previously unrated items reach the top-10 to more than four times
that of the strongest published baseline, while nine of the eleven
baselines never surface such an item at all. It remains behind one
recent graph-based hybrid on overall NDCG.

Future work: evaluation on hotel interaction logs once the platform
accumulates them; widening the search range for the item-side
constant; estimating the constants from the ratio of component error
variances rather than by grid search; and bringing the knowledge
pre-filter into the evaluated pipeline.

%=====================================================================
\begin{thebibliography}{00}

\bibitem{bernardi2015} L. Bernardi, J. Kamps, J. Kiseleva, and
M. J. I. Müller, ``The continuous cold start problem in e-commerce
recommender systems,'' in \textit{Proc. 2nd Workshop on New Trends
on Content-Based Recommender Systems}, 2015.

\bibitem{burke2002} R. Burke, ``Hybrid recommender systems: Survey
and experiments,'' \textit{User Modeling and User-Adapted
Interaction}, vol. 12, no. 4, pp. 331--370, 2002.

\bibitem{claypool1999} M. Claypool, A. Gokhale, T. Miranda,
P. Murnikov, D. Netes, and M. Sartin, ``Combining content-based and
collaborative filters in an online newspaper,'' in \textit{Proc. ACM
SIGIR Workshop on Recommender Systems}, 1999.

\bibitem{dacrema2019} M. Ferrari Dacrema, P. Cremonesi, and
D. Jannach, ``Are we really making much progress? A worrying analysis
of recent neural recommendation approaches,'' in \textit{Proc. 13th
ACM Conf. Recommender Systems}, 2019, pp. 101--109.

\bibitem{do2020} H.-Q. Do, T.-H. Le, and B. Yoon, ``Dynamic weighted
hybrid recommender systems,'' in \textit{Proc. 22nd Int. Conf.
Advanced Communication Technology (ICACT)}, 2020, pp. 644--650.

\bibitem{harper2015} F. M. Harper and J. A. Konstan, ``The MovieLens
datasets: History and context,'' \textit{ACM Trans. Interactive
Intelligent Systems}, vol. 5, no. 4, pp. 1--19, 2015.

\bibitem{herlocker2004} J. L. Herlocker, J. A. Konstan,
L. G. Terveen, and J. T. Riedl, ``Evaluating collaborative filtering
recommender systems,'' \textit{ACM Trans. Information Systems},
vol. 22, no. 1, pp. 5--53, 2004.

\bibitem{koren2008} Y. Koren, ``Factorization meets the
neighborhood: A multifaceted collaborative filtering model,'' in
\textit{Proc. 14th ACM SIGKDD Int. Conf. Knowledge Discovery and Data
Mining}, 2008, pp. 426--434.

\bibitem{krichene2020} W. Krichene and S. Rendle, ``On sampled
metrics for item recommendation,'' in \textit{Proc. 26th ACM SIGKDD
Int. Conf. Knowledge Discovery and Data Mining}, 2020,
pp. 1748--1757.

\bibitem{lops2011} P. Lops, M. de Gemmis, and G. Semeraro,
``Content-based recommender systems: State of the art and trends,''
in \textit{Recommender Systems Handbook}. Springer, 2011,
pp. 73--105.

\bibitem{melville2002} P. Melville, R. J. Mooney, and R. Nagarajan,
``Content-boosted collaborative filtering for improved
recommendations,'' in \textit{Proc. 18th National Conf. Artificial
Intelligence (AAAI)}, 2002, pp. 187--192.

\bibitem{ohwsf2025} M. Muhammad, Gunardi, Danardono, and D. Rosadi,
``Optimization hybrid weighted switching filtering (OHWSF) using SVD
and SVD++ for addressing data sparsity,'' \textit{Int. J. Advances in
Intelligent Informatics}, vol. 11, no. 3, p. 533, 2025.

\bibitem{pazzani1999} M. J. Pazzani, ``A framework for
collaborative, content-based and demographic filtering,''
\textit{Artificial Intelligence Review}, vol. 13, no. 5--6,
pp. 393--408, 1999.

\bibitem{schein2002} A. I. Schein, A. Popescul, L. H. Ungar, and
D. M. Pennock, ``Methods and metrics for cold-start
recommendations,'' in \textit{Proc. 25th ACM SIGIR Conf. Research and
Development in Information Retrieval}, 2002, pp. 253--260.

\bibitem{slarec2026} M. E. Çakır, Z. Çetinkaya, F. Horasan, and
A. H. Yurttakal, ``Adaptive weighting-based hybrid recommender
system: Self-learning adaptive recommendation (SL-ARec),''
\textit{PeerJ Computer Science}, vol. 12, p. e3856, 2026.

\bibitem{tarek2022} I. H. M. Tarek, F. U. Munna,
A. T. M. T. H. Mojumder, M. M. Rahman, M. S. Hossain, and
K. Andersson, ``A hybrid hotel recommendation using collaborative,
content based and knowledge based approach,'' in \textit{Intelligent
Computing \& Optimization}, Lecture Notes in Networks and Systems.
Springer, 2022, pp. 1049--1057.

\end{thebibliography}

\end{document}
